\documentclass[11pt,a4paper]{article}
% =====================================================================
%  Shared preamble for the Part V lecture notes (lectures 19-22).
%
%  These four lectures sit outside Geron, so the notes ARE the primary
%  source and are examined on the same terms as the textbook parts.
%  Everything here is chosen to make that readable on paper: one column,
%  generous measure, and the same six-colour palette as the slides so a
%  student moving between the two is not re-learning what red means.
%
%  Build with pdflatex (twice, for the toc and the refs):
%      cd notes && latexmk -pdf lecture-19.tex
%  or  make -C notes
% =====================================================================
\usepackage[T1]{fontenc}
\usepackage[utf8]{inputenc}
\usepackage{newtxtext,newtxmath}      % Times-like: prints well, reads well
\usepackage[scaled=0.86]{inconsolata} % code, at a size that sits beside Times
\usepackage{microtype}
% newtx defines \Bbbk, \openbox, \proof and \endproof; amssymb and amsthm
% then redefine them and abort. Freeing the four names before those two
% packages load is the standard fix, and it costs nothing here: we use
% amsthm only for \newtheorem.
\let\Bbbk\relax
\usepackage{amsmath,amssymb}
\let\openbox\relax
\let\proof\relax
\let\endproof\relax
\usepackage{amsthm}
\usepackage{booktabs}
\usepackage{enumitem}
\usepackage{xcolor}
\usepackage{tcolorbox}
\usepackage{titlesec}
\usepackage{graphicx}
\usepackage[hidelinks]{hyperref}
\tcbuselibrary{skins,breakable}

% ---- the course palette, from assets/css/custom.css -------------------
\definecolor{aimlprimary}{HTML}{0B3D62}   % deep blue  - structure
\definecolor{aimlaccent} {HTML}{C0392B}   % brick red  - failures
\definecolor{aimlsuccess}{HTML}{14663A}   % green      - repairs
\definecolor{aimlmath}   {HTML}{6C3483}   % purple     - the derivation
\definecolor{aimlmuted}  {HTML}{4B5563}
\definecolor{aimlrule}   {HTML}{B0BCC7}
\definecolor{aimlsoft}   {HTML}{F4F7F9}

\usepackage[a4paper,margin=32mm,top=30mm,bottom=30mm]{geometry}
\setlength{\parskip}{0.45\baselineskip}
\setlength{\parindent}{0pt}
\linespread{1.06}

% ---- headings --------------------------------------------------------
\titleformat{\section}{\Large\bfseries\color{aimlprimary}}{\thesection}{0.7em}{}
\titleformat{\subsection}{\large\bfseries\color{aimlprimary}}{\thesubsection}{0.7em}{}
\titleformat{\subsubsection}{\bfseries\color{aimlmuted}}{}{0em}{}
\titlespacing*{\section}{0pt}{1.6\baselineskip}{0.5\baselineskip}
\titlespacing*{\subsection}{0pt}{1.2\baselineskip}{0.35\baselineskip}

% ---- boxes -----------------------------------------------------------
% One box per role, and the role is the colour. A student who has seen the
% slides already knows what each one means before reading a word of it.
\newtcolorbox{keybox}[1][]{
  breakable, enhanced, colback=aimlsoft, colframe=aimlprimary,
  boxrule=0.4pt, left=8pt, right=8pt, top=6pt, bottom=6pt,
  fonttitle=\bfseries\small, coltitle=white, colbacktitle=aimlprimary,
  attach boxed title to top left={yshift=-2mm, xshift=4mm},
  boxed title style={boxrule=0pt, arc=1pt}, #1}

\newtcolorbox{failbox}[1][]{
  breakable, enhanced, colback=white, colframe=aimlaccent,
  boxrule=0.9pt, left=8pt, right=8pt, top=6pt, bottom=6pt,
  fonttitle=\bfseries\small, coltitle=white, colbacktitle=aimlaccent,
  attach boxed title to top left={yshift=-2mm, xshift=4mm},
  boxed title style={boxrule=0pt, arc=1pt}, #1}

\newtcolorbox{derivbox}[1][]{
  breakable, enhanced, colback=white, colframe=aimlmath,
  boxrule=0.6pt, left=8pt, right=8pt, top=6pt, bottom=6pt,
  fonttitle=\bfseries\small, coltitle=white, colbacktitle=aimlmath,
  attach boxed title to top left={yshift=-2mm, xshift=4mm},
  boxed title style={boxrule=0pt, arc=1pt}, #1}

% An exercise. Deliberately the same shape as the notebook's `try` field:
% one change, and what should happen to the output.
\newtcolorbox{trybox}[1][]{
  breakable, enhanced, colback=white, colframe=aimlsuccess,
  boxrule=0.5pt, left=8pt, right=8pt, top=6pt, bottom=6pt,
  fonttitle=\bfseries\small, coltitle=white, colbacktitle=aimlsuccess,
  attach boxed title to top left={yshift=-2mm, xshift=4mm},
  boxed title style={boxrule=0pt, arc=1pt}, #1}

% ---- small semantics -------------------------------------------------
\newcommand{\scope}[1]{\textcolor{aimlmuted}{\small #1}}
\newcommand{\term}[1]{\textbf{#1}}
\newcommand{\code}[1]{\texttt{\small #1}}
\newcommand{\lecture}[1]{Lecture~#1}

\newtheorem{definition}{Definition}[section]
\newtheorem{proposition}[definition]{Proposition}

% ---- title block -----------------------------------------------------
% \notestitle{19}{Information retrieval: the lexical foundation}{SciFact (BEIR)}
\newcommand{\notestitle}[3]{%
  \thispagestyle{empty}
  \begin{flushleft}
    {\color{aimlmuted}\small\sffamily
     APPLICATIONS OF MACHINE LEARNING \quad$\cdot$\quad
     BSc Mathematics of Artificial Intelligence}\\[2mm]
    {\color{aimlrule}\rule{\textwidth}{0.8pt}}\\[4mm]
    {\color{aimlmuted}\large Lecture #1 \quad$\cdot$\quad Part V \quad$\cdot$\quad
     Extended notes}\\[3mm]
    {\Huge\bfseries\color{aimlprimary}\baselineskip=1.15\baselineskip #2\par}\vspace{4mm}
    {\color{aimlmuted}\large Dataset: #3}\\[3mm]
    {\color{aimlrule}\rule{\textwidth}{0.8pt}}
  \end{flushleft}
  \vspace{2mm}
  \begin{keybox}[title={Why these notes exist}]
    Lectures 19--22 sit outside G\'eron, the course's primary text. For those
    four lectures \emph{these notes are the primary source}, and they are
    examined on exactly the same terms as the textbook parts. Everything in
    the slide deck for this lecture is here, with the arguments written out at
    length rather than compressed onto a slide; the numbers are the ones the
    accompanying notebook prints, and every one of them is a measurement on
    the corpus named above rather than a figure quoted from a paper.
  \end{keybox}
  \vspace{1mm}
  {\small\tableofcontents}
  \vspace{2mm}
  \noindent{\color{aimlrule}\rule{\textwidth}{0.4pt}}
  \vspace{1mm}
}


\begin{document}
\notestitle{21}{Recommender systems: from ratings to factors}{MovieLens 1M}

\section{The problem, and the data}

6{,}040 users. 3{,}706 films. 1{,}000{,}209 ratings on a 1--5 scale. Two
questions, and they are not the same question.

\begin{itemize}[leftmargin=*,nosep]
  \item \term{Rating prediction} --- what would this user give this film?
        Scored by RMSE. That is this lecture.
  \item \term{Top-$N$ recommendation} --- which ten films should we show?
        Scored by the ranking metrics of \lecture{19}. That is \lecture{22}.
\end{itemize}

They come apart, which is the subject of the next lecture: a model that
predicts every rating within half a star can still recommend badly, because the
films a user would rate highly are not the films they have not seen.

\subsection{Why MovieLens, and what it is not}

It has \emph{explicit ratings}, which most systems do not have --- but without
them, rating prediction cannot be taught at all. It has \emph{timestamps},
which is why this release and not the smaller one: half of \lecture{22} is the
difference between splitting a history at random and splitting it in time. And
a million ratings is large enough that the results are not noise and small
enough to factorise on a CPU in seconds.

\begin{keybox}[title={And what it is not}]
A research dataset from a system that \emph{asked people to rate films}. It is
not a log of a production recommender, its users are self-selected, and its
sparsity pattern is unlike a real catalogue's. Everything measured here
transfers as a \emph{method}, not as a number.
\end{keybox}

\subsection{The matrix, and what is not in it}

\begin{center}
\small
\begin{tabular}{lr}
\toprule
Users $\times$ films        & 22.4 million cells \\
Ratings present             & 1{,}000{,}209 \\
Density                     & 4.5\% \\
Ratings per user, median    & 96 \\
Ratings per film, median    & 124 \\
\bottomrule
\end{tabular}
\end{center}

95.5\% of the matrix is missing, and this is the central fact of the lecture.

\begin{keybox}[title={Missing is not zero}]
A cell is empty because the user never saw the film, or saw it and did not rate
it, or was never shown it by whatever system was running at the time. Every one
of those is a different meaning, and none of them is ``disliked''.
\end{keybox}

\subsubsection{Not missing at random}

Worse: the pattern of what is missing is \emph{informative}. People rate films
they chose to watch, and they chose them because they expected to like them. A
horror film unrated by a user who hates horror is missing \emph{because} they
hate horror --- which is a rating of a kind, and it is not in the data. Every
model in this lecture trains on observed entries and is evaluated on observed
entries, so none of them ever sees this.

That is the same shape as \lecture{19}'s pooling: the data records what someone
chose to look at, and the choice was not random. In retrieval it made recall a
lower bound; here it biases the training set itself.

\subsection{What ``a user'' means in the model}

$p_u$ is a vector attached to an \emph{account}, and an account is not a
person.

\begin{itemize}[leftmargin=*,nosep]
  \item A household shares one account, so $p_u$ averages several tastes into
        one point --- and the average of two tastes is often nobody's.
  \item A person's taste changes; $p_u$ is a single vector with no time in it,
        and MovieLens spans years.
  \item The same person on two devices is two users, with two independent
        vectors estimated from half the data each.
\end{itemize}

Nothing in the objective represents any of this. A model whose unit of analysis
is wrong will be confidently wrong, and no amount of held-out RMSE reveals it.

\subsection{What people actually give}

\begin{center}
\small
\begin{tabular}{lccccc}
\toprule
Rating & 1 & 2 & 3 & 4 & 5 \\
Share  & 5.6\% & 10.8\% & 26.1\% & 34.9\% & 22.6\% \\
\bottomrule
\end{tabular}
\end{center}

Mean 3.58, heavily skewed upward: 57.5\% of ratings are 4 or 5. People mostly
rate things they liked, which means the scale is not being used as a scale ---
it is being used as ``good'', ``very good'', and occasionally ``I want you to
know this was bad''.

\subsection{Why RMSE, and its quiet assumption}

\[
  \text{RMSE} = \sqrt{\frac{1}{|\Omega_{\text{test}}|}
  \sum_{(u,i)\in\Omega_{\text{test}}} \big(r_{ui} - \hat r_{ui}\big)^2}
\]

Squared error, so a two-star miss costs four times a one-star miss. That is
defensible and worth stating: it says being badly wrong occasionally is worse
than being slightly wrong often.

The quiet assumption is bigger. RMSE averages over the \emph{observed} test
ratings --- films these users chose to watch and rate. It says nothing about
the films they did not, which is precisely the set a recommender operates on.
\lecture{22} replaces it for that reason. Today it is the right metric for
``what would they rate it'' and the wrong one for ``what should we show''.

\subsection{The long tail}

\begin{center}
\small
\begin{tabular}{lr}
\toprule
\textbf{Top films by rating count} & \textbf{Share of all ratings} \\
\midrule
top 1\%  & 8.7\% \\
top 5\%  & 28.3\% \\
top 10\% & 44.4\% \\
top 20\% & 65.2\% \\
\midrule
Films with fewer than 20 ratings & 17.9\% \\
\bottomrule
\end{tabular}
\end{center}

Two consequences to keep. First, a model can score well by learning popularity
and nothing else --- \lecture{22} measures exactly that. Second, the films
where a recommendation would be \emph{valuable} are the ones with least data.

\begin{keybox}[title={The tail is where the value is}]
A recommendation is useful to the extent the user would not have found the item
anyway. The top 1\% of films hold 8.7\% of ratings: recommending those is easy,
scores well, and tells the user nothing they did not know. And 17.9\% of films
have under twenty ratings --- exactly where every model here has least to work
with.

Nothing measured in this lecture rewards usefulness. RMSE over held-out
ratings, and even NDCG over held-out interactions, give full credit for a
correct popular recommendation and the same credit for a correct obscure one.
\end{keybox}

\section{Baselines: how much of a rating is the user, and how much the film}

\subsection{The split, stated first}

90/10 on the ratings, at random, one seed, fixed before any model was fitted:
900{,}189 ratings to fit and 100{,}020 to test. Every number in this lecture
uses that split.

\begin{failbox}[title={A random split of ratings is not innocent}]
It puts a user's \emph{future} ratings in the training set and their past in
the test set, which no deployed system ever gets. It is the standard protocol
for rating prediction, it is convenient, and it is optimistic. \lecture{22}
measures exactly how optimistic.
\end{failbox}

\subsection{The baselines}

\begin{center}
\small
\begin{tabular}{lr}
\toprule
\textbf{Predictor} & \textbf{RMSE} \\
\midrule
The global mean, for everyone & 1.1185 \\
This user's mean rating       & 1.0367 \\
This film's mean rating       & 0.9826 \\
Both, as biases               & 0.9109 \\
\bottomrule
\end{tabular}
\end{center}

A single number for everybody scores 1.1185. Knowing which film takes it to
0.9826; knowing who \emph{and} which film takes it to 0.9109. No interaction
between the two has been modelled yet --- nothing so far knows that
\emph{this} user likes \emph{this kind} of film.

\subsection{The bias model, written out}

\[ \hat r_{ui} = \mu + b_u + b_i \]

Fitted by alternating closed forms, each regularised by the count, because a
user with three ratings should not be given a large offset:

\[
  b_u = \frac{\sum_{i \in I_u}\big(r_{ui} - \mu - b_i\big)}{|I_u| + \lambda}.
\]

That $\lambda$ in the denominator is shrinkage --- the ridge penalty of
\lecture{5} arriving in a new place: with few observations, pull the estimate
toward zero.

\begin{center}
\small
\begin{tabular}{lrr}
\toprule
 & \textbf{Lowest} & \textbf{Highest} \\
\midrule
User offset $b_u$ & $-2.2156$ & $+1.4690$ \\
Film offset $b_i$ & $-1.8545$ & $+1.0546$ \\
\bottomrule
\end{tabular}
\end{center}

The harshest user rates nearly 2.2 stars below the mean on everything and the
most generous 1.5 above. That spread is larger than most of what the rest of
this lecture will add.

\begin{keybox}[title={Why this matters more than it looks}]
A recommender that ignores biases spends its capacity re-learning that some
users are generous and some films are popular --- effects with nothing personal
in them. Modelling them explicitly and cheaply frees the interesting part of
the model for the interesting part of the problem.
\end{keybox}

\subsection{Cold start, named now}

A user with no ratings has no $p_u$ to solve for; a film with no ratings has no
$q_i$. The factorisation has nothing to say about either, and the failure is
total rather than gradual. A new user falls back to popularity, or you ask for
a few ratings up front. A new item falls back to content: genre, cast,
description, an embedding of the synopsis. A new \emph{system} has neither,
which is why every recommender starts life as a popularity list --- and why a
real one is always a hybrid of a factorisation and something that works without
history.

\section{Derivation: matrix factorisation, and the SVD it is not}

\subsection{What we are given}

A matrix $R \in \mathbb{R}^{m\times n}$ with $m = 6{,}040$ users and
$n = 3{,}706$ films, of which we observe a set $\Omega$ of entries:
\[ |\Omega| = 1{,}000{,}209, \qquad \frac{|\Omega|}{mn} = 4.5\%. \]

$R$ is \emph{not} a sparse matrix in the sense of \lecture{19}, where the zeros
were real zeros. Here the unobserved entries have no value at all --- not zero,
not the mean, not anything. Every difficulty in this derivation comes from that
one sentence, and every wrong method comes from forgetting it.

\subsection{The low-rank hypothesis}

Assume tastes are simpler than the data: that there are $d$ underlying factors
--- how much a film is a comedy, how dark, how old --- and that a rating is how
much the user wants each factor times how much the film has of it.
\[ \hat r_{ui} = p_u \cdot q_i, \qquad p_u, q_i \in \mathbb{R}^d, \]
equivalently $\hat R = PQ^{\top}$ with $\operatorname{rank} \le d$. The
parameter count falls from $mn$ to $d(m+n)$, which at $d = 16$ is about a
hundredth of the matrix.

It is a \emph{hypothesis}, not a fact. Its justification is that it predicts
held-out ratings better than the alternatives, and we check that rather than
assert it.

\subsection{Lecture 8 already solved this --- or did it}

The best rank-$d$ approximation of a matrix in the Frobenius norm is its
truncated SVD. Eckart--Young:
\[
  \min_{\operatorname{rank}(X)\le d} \|R - X\|_F
  \;=\; \big\|R - U_d\Sigma_d V_d^{\top}\big\|_F .
\]
So the answer looks as though it is already in our hands: take the SVD of the
ratings matrix, keep $d$ singular values, read off $P$ and $Q$.

\begin{failbox}[title={One condition, and it is not met}]
The Frobenius norm sums over \emph{every} entry of $R$. To evaluate it --- let
alone minimise it --- every entry must have a value, and 95.5\% of ours do not.
\end{failbox}

Be precise about what the theorem says, because the derivation turns on it. It
is a statement about \emph{one given matrix} $R$; the norm sums over every
entry with equal weight; and it says nothing about entries $R$ does not have,
because the question is not defined. The theorem is not wrong. It is being
misused. Filling the matrix does not extend the theorem to our problem --- it
\emph{changes our problem} into one the theorem covers, and the two problems
have different answers.

\subsection{So fill them in?}

Two obvious choices, both actually tried, both measured on the same held-out
ratings.

\begin{center}
\small
\begin{tabular}{lr}
\toprule
\textbf{Rank-32 SVD of\ldots} & \textbf{RMSE} \\
\midrule
the matrix filled with zeros           & 2.3274 \\
the matrix filled with the global mean & 0.9920 \\
\midrule
predicting the global mean for everyone & 1.1185 \\
the bias model                          & 0.9109 \\
\bottomrule
\end{tabular}
\end{center}

Filling with zeros scores 2.3274 --- more than twice the error of predicting
one number for everybody, \emph{from the mathematically optimal rank-32
approximation}.

Ratings run from 1 to 5, so a zero is not a low rating: it is off the scale.
The approximation is dragged toward zero everywhere and the predictions for
observed pairs come out far too low. Filling with the mean at least puts the
fill inside the range, so the damage is bounded: 0.9920, merely worse than the
bias model rather than catastrophic. Both are the same error --- one is simply
more visible, and the more visible error is the safer one to make.

\begin{keybox}[title={The transferable error}]
Imputing missing values to make a method applicable makes the method answer a
different question. It is the same mistake as filling a missing feature with a
column mean and then reporting the model as if the data had been complete.
\end{keybox}

\subsection{Change the objective, not the data}

The problem was never the SVD. It was that we asked for a fit to \emph{every}
entry. So ask only for what we can observe:

\begin{derivbox}[title={The factorisation objective}]
\[
  \min_{P,Q} \sum_{(u,i)\in\Omega} \big(r_{ui} - p_u\cdot q_i\big)^2
  \;+\; \lambda\big(\|P\|_F^2 + \|Q\|_F^2\big)
\]
\begin{itemize}[leftmargin=*,nosep]
  \item $\Omega$ --- the observed entries, \emph{and only those}. This is the
        whole departure from the SVD.
  \item $p_u\cdot q_i$ --- a rank-$d$ model written entrywise rather than as a
        matrix product, because the matrix product is over cells we do not
        have.
  \item $\lambda$ --- without which a user with three ratings gets a perfect
        fit and no information.
\end{itemize}
Three symbols, three decisions. If you can reconstruct why each is there, you
have the derivation.
\end{derivbox}

The unobserved entries are now \emph{absent} from the objective rather than
assigned a value. Nothing is imputed and nothing is assumed about what a user
would have said. And in doing so we have left the world where the SVD applies.

\subsection{What was given up}

\begin{center}
\small
\begin{tabular}{p{0.13\textwidth}p{0.33\textwidth}p{0.4\textwidth}}
\toprule
 & \textbf{Truncated SVD} & \textbf{This objective} \\
\midrule
Solution   & closed form & iterative \\
Uniqueness & unique (up to sign and ties) & not unique --- $PM$, $QM^{-\top}$
score identically \\
Factors    & orthogonal, ordered by variance & neither \\
Optimum    & global & local --- the problem is non-convex in $(P,Q)$ jointly \\
\bottomrule
\end{tabular}
\end{center}

All four losses are real, and the last two matter in practice: you cannot
interpret factor 3 the way you interpret principal component 3, because a
different run gives different, equally good factors.

\subsection{Regularisation is not optional}

A user with three ratings and $d = 16$ has more parameters than observations,
so the objective can be driven to zero on that user while learning nothing. The
penalty is ridge regression again, from \lecture{5}, doing the same job: with
few observations, prefer the smaller coefficients. It also has a second job
here --- it breaks the scaling degeneracy, since $P \to 2P$, $Q \to Q/2$ is no
longer free.

\subsection{Hold one side fixed}

The objective is not convex in $P$ and $Q$ together. But fix $Q$, and every
term involving $p_u$ is a quadratic in $p_u$ alone:
\[
  \min_{p_u} \sum_{i \in I_u} \big(r_{ui} - p_u\cdot q_i\big)^2
  + \lambda\|p_u\|^2 ,
\]
which is a \emph{ridge regression}, with the films this user rated as the rows
of the design matrix and their ratings as the target. It is separate for each
user: no user's solution depends on another's.

\begin{derivbox}[title={The ALS half-step}]
Write $Q_u$ for the $|I_u| \times d$ matrix of the films user $u$ rated and
$r_u$ for their ratings. Setting the gradient to zero,
\[
  p_u = \big(Q_u^{\top}Q_u + \lambda I\big)^{-1} Q_u^{\top} r_u .
\]
The normal equation of \lecture{2} with the ridge term of \lecture{5}. The
system is $d\times d$ --- $16\times16$ here --- regardless of how many films
the user rated.

Note where the missing entries went: $Q_u$ contains \emph{only the films this
user rated}. Nothing was imputed; the unobserved entries simply never enter the
system.
\end{derivbox}

By symmetry, fixing $P$ makes each $q_i$ a ridge regression over the users who
rated film $i$. So: fix $Q$ and solve for every $p_u$; fix $P$ and solve for
every $q_i$; repeat.

\begin{proposition}
Each half-step solves its subproblem exactly, so neither can increase the
objective. The objective is bounded below by zero, so the sequence of objective
values converges.
\end{proposition}

It converges to a \emph{local} minimum: monotone descent is not global
optimality, and a different initialisation gives a different answer. That is
the price of changing the objective, and it was paid knowingly.

\subsubsection{What a sweep costs}

\begin{center}
\small
\begin{tabular}{lr}
\toprule
Systems solved per sweep & 6{,}040 users $+$ 3{,}706 films \\
Size of each system      & $d \times d$ \\
At $d = 16$              & $16\times16$ \\
Ratings gathered per sweep & 900{,}189, twice \\
\bottomrule
\end{tabular}
\end{center}

Each system is tiny and independent of the others, so a sweep is embarrassingly
parallel and its cost is dominated by \emph{gathering} each user's rows rather
than by the linear algebra --- which is why the notebook groups the ratings
once with a single \code{argsort} instead of taking a boolean mask per user.
The same arithmetic, and the difference between seconds and minutes.

\subsection{Or descend, one rating at a time}

For each observed rating, take the error and step both vectors:
\[
  e_{ui} = r_{ui} - \hat r_{ui}, \qquad
  p_u \mathrel{+}= \eta\,(e_{ui}q_i - \lambda p_u), \qquad
  q_i \mathrel{+}= \eta\,(e_{ui}p_u - \lambda q_i).
\]

ALS parallelises across users and each step is exact, but needs a $d\times d$
solve per user. SGD touches one rating at a time, so it suits implicit data and
streams, but needs a learning rate.

\begin{failbox}[title={Copy $p_u$ before you update it}]
Using the already-updated $p_u$ in the $q_i$ step is a different algorithm. It
still converges, and nothing warns you.
\end{failbox}

Same objective, different route. At $d = 32$, ALS reaches 0.8733 and SGD
0.8615. Fifteen epochs of per-rating updates against twelve ALS sweeps: two
routes to one objective, arriving in comparable places.

\subsection{Put the biases back}

The factorisation models the \emph{interaction}. The parts that are not
interactions --- a generous user, a popular film --- should not be spent on it:
\[ \hat r_{ui} = \mu + b_u + b_i + p_u\cdot q_i . \]

\begin{center}
\small
\begin{tabular}{lr}
\toprule
Bias model alone ($d = 0$)                & 0.9109 \\
Factorisation, $d = 16$, with biases      & 0.8587 \\
Factorisation, $d = 16$, without biases   & 0.8559 \\
\bottomrule
\end{tabular}
\end{center}

The biases alone close 79.9\% of the distance from the global mean to the full
model. \emph{Most of a rating is not personal.}

\begin{keybox}[title={An honest wrinkle}]
At $d = 16$ the model \emph{without} explicit biases scores slightly better
than the one with them. The textbook claim is that biases help; here they do
not, and the reason is visible in the algebra --- a factor dimension can hold a
constant, so a rank-16 factorisation can represent the biases itself, using two
of its sixteen dimensions to do it.

What is still true: biases are worth having because they are cheap, stable and
interpretable, they help most at low rank, and they let the factors be about
interaction. Not because they always lower RMSE --- and the difference here,
0.0028, is small enough that reporting it as a finding either way would be
overreaching.
\end{keybox}

\subsection{How many factors?}

\begin{center}
\small
\begin{tabular}{lccccccc}
\toprule
$d$   & 1 & 2 & 4 & 8 & 16 & 32 & 64 \\
RMSE  & 0.8912 & 0.8836 & 0.8705 & 0.8602 & 0.8587 & 0.8733 & 0.9006 \\
\bottomrule
\end{tabular}
\end{center}

Best at $d = 16$, and rising again by $d = 64$. A U-curve, which is
\lecture{5} exactly: more capacity fits the training ratings better and the
held-out ratings worse. The rank is a capacity parameter and it is chosen the
way every capacity parameter in this course is chosen --- by measurement on
data the model did not see.

Even $d = 1$ scores 0.8912, already better than the bias model. One dimension
of taste is worth something.

\subsection{What the factors are not}

\begin{itemize}[leftmargin=*,nosep]
  \item They are not principal components: not orthogonal, not ordered, not
        unique.
  \item They are not genres. A run may put something genre-like in a dimension,
        and the next run will not put it in the same one.
  \item They are not interpretable individually.
\end{itemize}

\begin{derivbox}[title={The rotation nobody can see}]
For any invertible $M$,
\[ (PM)(QM^{-\top})^{\top} = PMM^{-1}Q^{\top} = PQ^{\top} . \]
Every prediction is identical, so the objective cannot distinguish $P$ from
$PM$. The solution is a whole family, and an optimiser returns whichever member
its initialisation led to. The \emph{predictions} are determined --- that is
what we wanted. The \emph{factors} are not, and no amount of staring at them
fixes that.

The regularisation narrows the family: it penalises $\|P\|$, so $M$ can no
longer be an arbitrary scaling. Rotations of an isotropic penalty remain free.
\end{derivbox}

Say this out loud when someone shows you a factor plot. Reading meaning into an
individual latent dimension is the recommender-systems version of reading a
coefficient of a collinear regression: the model is not claiming what the plot
claims.

\section{Neighbourhoods, and feedback nobody gave}

\subsection{The other family}

Before factorisation there was a simpler idea, and it is still deployed:
predict from similar things. \term{User-based} finds users who rated as you did
and averages their ratings for the film in question; \term{item-based} finds
films rated as this film was and averages this user's ratings of them.

Item-based is what gets deployed, for a reason that is operational rather than
statistical: there are usually far fewer items than users, item--item
similarities change slowly, and they can be computed offline and cached.

\subsection{Item--item similarity}

\[
  s_{ij} = \frac{\tilde r_i \cdot \tilde r_j}{\|\tilde r_i\|\,\|\tilde r_j\|}
\]
--- the cosine between two films' rating columns, where the tilde matters: the
columns are \emph{bias-adjusted} first,
$\tilde r_{ui} = r_{ui} - \mu - b_u - b_i$. Without that, every popular film is
similar to every other popular film, because both columns are full of fours.

That is the same trap as an unnormalised dot product in \lecture{20} and as
term frequency without idf in \lecture{19}: subtract what is true of everything
before measuring what is specific.

\[
  \hat r_{ui} = \mu + b_u + b_i +
  \frac{\sum_{j\in N_k(i)} s_{ij}\,\tilde r_{uj}}{\sum_{j\in N_k(i)} s_{ij}}
\]

\begin{center}
\small
\begin{tabular}{lr}
\toprule
\textbf{Method} & \textbf{RMSE} \\
\midrule
item--item, $k = 10$        & 0.9376 \\
item--item, $k = 20$        & 0.9171 \\
item--item, $k = 50$        & 0.8867 \\
factorisation, $d = 16$     & 0.8587 \\
\bottomrule
\end{tabular}
\end{center}

The best neighbourhood model does not reach the factorisation, and still beats
the bias model.

\subsection{Why factorisation wins, and where it does not}

\begin{center}
\small
\begin{tabular}{p{0.16\textwidth}p{0.33\textwidth}p{0.37\textwidth}}
\toprule
 & \textbf{Neighbourhood} & \textbf{Factorisation} \\
\midrule
Uses         & direct co-rating evidence & a global low-rank structure \\
Sparse items & few co-ratings, poor estimate & borrows strength from
everything \\
Explainable  & ``because you liked X'' & a dot product of latent vectors \\
New rating   & update a cached similarity & refit, or fold in approximately \\
\bottomrule
\end{tabular}
\end{center}

Row three is why neighbourhood methods survive: a recommendation a user can be
told the reason for is worth points a metric does not award. And in practice
the two are combined --- exactly as BM25 and dense retrieval were combined in
\lecture{20}, and for the same reason.

\subsubsection{Where a neighbourhood method runs out}

Two films are similar only if enough users rated \emph{both}. In the long tail
there are no such users: 17.9\% of films have fewer than twenty ratings at all,
for two such films the co-rating count is typically zero, and the cosine of two
nearly empty columns is noise.

\begin{keybox}[title={Borrowing strength}]
A factorisation has no such problem: a tail film's $q_i$ is estimated from its
few ratings \emph{and from where those raters sit in a space learned from
everything else}. That is the real advantage of the low-rank model, and it is
the same argument as pooling in a hierarchical model --- a parameter estimated
from three observations is rescued by a structure estimated from a million.
\end{keybox}

\subsection{Explicit ratings are the easy case}

\begin{center}
\small
\begin{tabular}{p{0.42\textwidth}p{0.48\textwidth}}
\toprule
\textbf{Explicit} & \textbf{Implicit} \\
\midrule
a rating, deliberately given & a click, a play, a purchase, a dwell time \\
scarce --- most users rate nothing & abundant --- every interaction is one \\
a scale, with a middle & presence or absence \\
a stated preference & a revealed one, plus the interface \\
\bottomrule
\end{tabular}
\end{center}

MovieLens is explicit, which is why it can be used to teach rating prediction
at all. It is not representative: your first recommender will almost certainly
have no ratings in it.

Binarise our ratings at 4 or above and call that an interaction:

\begin{center}
\small
\begin{tabular}{lr}
\toprule
Explicit ratings           & 1{,}000{,}209 \\
Of which positive ($\ge 4$) & 575{,}281 \\
As a fraction              & 57.5\% \\
Cells with no interaction  & 22.4 million \\
\bottomrule
\end{tabular}
\end{center}

In a real system the first row would be zero and the second would be a click
log ten times larger. The proportions are what matter: a few hundred thousand
positives against tens of millions of cells whose meaning is unknown. This
binarised view is what \lecture{22} uses throughout, which is why it is defined
here.

\subsection{Implicit feedback has no negatives}

A user watched film A. What does that say about film B, which they did not
watch? They disliked it --- a true negative. Or they have not seen it yet ---
possibly the best recommendation available. Or they were never shown it, and
that was the old system's choice rather than theirs.

\begin{failbox}[title={The whole difficulty}]
In explicit data the missing entries are excluded from the objective, and we
saw exactly what happens when they are not. In implicit data the \emph{only
thing that exists} is the positives, so an objective over observed entries
alone is trivially satisfied by predicting ``yes'' for everything.
\end{failbox}

Two ways out.

\begin{enumerate}[leftmargin=*]
  \item \term{Weighted least squares over everything.} Treat every unobserved
        cell as a zero with \emph{low confidence}, and observed ones as a one
        with confidence rising in the number of interactions. All 22.4 million
        cells enter the objective, with weights. This keeps the closed form ---
        the weighting is arranged so that ALS still works --- and needs a
        confidence function you have to justify.
  \item \term{Pairwise ranking.} Never predict a value at all; require only
        that an observed item outranks an unobserved one for that user. This
        changes the objective into a ranking objective, which is what we
        actually want, and gives up the closed form.
\end{enumerate}

\subsubsection{BPR: rank the pair, not the item}

\begin{derivbox}[title={Bayesian personalised ranking}]
\[
  \max \sum_{(u,i,j)} \log \sigma\!\big(\hat r_{ui} - \hat r_{uj}\big)
  \;-\; \lambda\|\Theta\|^2 ,
\]
with $i$ observed for $u$ and $j$ sampled from the unobserved.
\begin{itemize}[leftmargin=*,nosep]
  \item Only the \emph{difference} is modelled, so no absolute score is ever
        claimed --- and none was ever observed.
  \item $j$ is \emph{sampled}, not enumerated, which is what makes it
        affordable.
  \item The logistic of a score difference is the pairwise logistic loss of
        \lecture{4}, and the sampling is the negative sampling of \lecture{20}.
\end{itemize}
\end{derivbox}

And it inherits \lecture{20}'s trap exactly: a sampled ``negative'' may be an
item the user would love. \emph{Being unobserved is not evidence of dislike.}

\section{Where we are}

\begin{center}
\small
\begin{tabular}{lrl}
\toprule
\textbf{Model} & \textbf{RMSE} & \textbf{Knows} \\
\midrule
Global mean              & 1.1185 & nothing \\
Film mean                & 0.9826 & which film \\
Bias model               & 0.9109 & who, and which film \\
SVD, mean-filled         & 0.9920 & an imputation it invented \\
Item--item, $k = 50$     & 0.8867 & co-rating evidence \\
Factorisation, $d = 16$  & 0.8587 & a low-rank structure \\
\bottomrule
\end{tabular}
\end{center}

Every row measured on the same held-out 100{,}020 ratings, with the same split
and the same seed.

\subsection{The distance travelled, in one number}

\begin{center}
\small
\begin{tabular}{lrr}
\toprule
 & \textbf{RMSE} & \textbf{Distance closed} \\
\midrule
Global mean             & 1.1185 & 0\% \\
Bias model              & 0.9109 & 79.9\% \\
Item--item, $k = 50$    & 0.8867 & 89.2\% \\
Factorisation, $d = 16$ & 0.8587 & 100\% \\
\bottomrule
\end{tabular}
\end{center}

Two arithmetic means and a shrinkage parameter get you most of the way; the
derivation, the alternating solves and sixteen latent dimensions get the rest.
That ratio is not an argument against the model. It is an argument for always
fitting the baseline first, so that you know which part of your result is the
idea and which part is arithmetic.

\subsection{What a deployed recommender actually is}

\begin{enumerate}[leftmargin=*,nosep]
  \item \term{Candidates}: catalogue $\rightarrow$ hundreds. Factorisation,
        popularity, co-visitation, recency --- several sources, unioned.
  \item \term{Rank}: hundreds $\rightarrow$ tens. A richer model with features
        the first stage cannot afford.
  \item \term{Policy}: diversity, freshness, business rules, exploration.
\end{enumerate}

The same three stages as \lecture{19}'s search pipeline, for the same reason:
the accurate model does not scale to the catalogue. What we built today is
stage 1. Stage 3 has no equation in it and decides more of what a user sees
than stages 1 and 2 together.

\begin{trybox}[title={Exercise 21.1}]
Set the regularisation to zero and re-run the rank sweep. Watch where the
U-curve moves and which rank stops being the best one. Then reproduce the
zero-filled SVD figure yourself --- it is the most surprising number in Part~V
and it takes four lines.
\end{trybox}

\section{Reference}

\subsection{Five questions for a recommender result}

\begin{enumerate}[leftmargin=*]
  \item \term{Rating prediction or ranking?} RMSE and NDCG answer different
        questions and are not comparable.
  \item \term{What is the baseline?} The bias model is four lines and is often
        close.
  \item \term{How was it split?} Randomly over ratings, or in time?
  \item \term{What happened to the missing entries?} Excluded, imputed, or
        weighted --- and if imputed, with what.
  \item \term{Explicit or implicit}, and if implicit, where did the negatives
        come from?
\end{enumerate}

Question three is the one \lecture{22} turns into a measurement.

\subsection{Where students lose marks}

\begin{itemize}[leftmargin=*,nosep]
  \item Saying ``matrix factorisation is the SVD''. It is not, and being able
        to say why in one sentence is the most-asked examination question in
        this part.
  \item Writing the objective as a sum over all $(u,i)$ rather than over
        $\Omega$.
  \item Claiming ALS finds the global optimum. Each half-step is exact; the
        joint problem is not convex.
  \item Interpreting a latent factor as a genre.
  \item Treating an unobserved entry as a negative in implicit data --- the
        error the whole last section is about.
\end{itemize}

\subsection{Vocabulary}

\begin{center}
\small
\begin{tabular}{ll}
\toprule
Explicit / implicit & a stated rating, against an observed interaction \\
Cold start          & a user or item with no history to factorise \\
Bias term           & the part of a rating that is not an interaction \\
Latent factor       & a coordinate of $p_u$ or $q_i$; individually meaningless \\
ALS                 & alternating ridge regressions, one side at a time \\
BPR                 & a pairwise ranking objective for implicit data \\
Long tail           & the majority of items, holding a minority of the interactions \\
\bottomrule
\end{tabular}
\end{center}

\subsection{Scope}

\term{Examinable.} The recommendation problem and how its data differs;
missing-not-at-random; the derivation in full --- the low-rank model, why the
SVD cannot be taken, the observed-entry objective, the ALS half-step and its
closed form, the SGD update, biases; the rank as a capacity parameter;
item--item neighbourhoods; why implicit feedback has no negatives, and BPR's
response.

\term{Not examinable --- engineering.} Recommender libraries, serving
infrastructure, incremental refitting.

\term{Beyond the syllabus.} SVD++ and time-aware factorisation, the full
weighted-ALS derivation for implicit data, bandits for exploration.

\subsection{Reading}

\begin{itemize}[leftmargin=*,nosep]
  \item These notes are the primary source.
  \item Koren, Bell and Volinsky, \emph{Matrix Factorization Techniques for
        Recommender Systems} --- the Netflix-era summary, and still the
        clearest account of biases.
  \item Hu, Koren and Volinsky, \emph{Collaborative Filtering for Implicit
        Feedback Datasets} --- the weighted-ALS method.
  \item Rendle et al., \emph{BPR} --- the pairwise objective.
\end{itemize}

\subsection{Summary}

\begin{keybox}[title={One line to keep}]
The missing entries \emph{are} the data. Impute them and you fit your own
assumption.
\end{keybox}

Recommendation is retrieval with the query replaced by a history, and the
labels are a by-product of use rather than a measurement. 95.5\% of the matrix
is missing, and missing is not zero: filling it and taking the SVD scores
2.3274, worse than predicting one number for everybody. Restrict the objective
to observed entries and you leave the SVD behind --- no closed form, no
uniqueness, no interpretable factors, and a model that works. ALS is a ridge
regression per user; SGD is one update per rating; same objective, different
route, comparable answers. And most of a rating is not personal: the biases
alone reach 0.9109 against the full model's 0.8587.

\end{document}
